\scanpage{23}
Alors on sait que l'on peut prendre les morphismes \inserted{de transition} $Z_j\unsure{\mathrel{\text{---}}|\longrightarrow} Z_i$ des \unsure{imm|imm.} fermées, et il suffit de prouver qu'alors les $f_*(Z_j)\mathrel{\text{---}} f_*(Z_i)$ sont aussi des imm. fermées, car à la limite on aura que $f_*(Z)\hookrightarrow f_*(Z_i)$ par imm. fermée, donc $f_*(Z)$ sera de t.f. sur $S$.

On est donc ramené à prouver en \unsure{m\textsuperscript{e}|même} temps le

\uline{Corollaire 20}. Si $Z\longrightarrow Z'$ est une imm. fermée de $X$-schémas affines, alors $f_*(Z)\longrightarrow f_*(Z')$ est une imm. fermée.

Pour prouver ceci, on est \unsure{également|aisément}, par dévissage standard, ramené au cas où $Z,Z'$ de présentation finie.

Ayant des situations de présentation finie, on est ramené au cas $S$ de type fini sur $\operatorname{Spec}\mathbb Z$, donc noeth.

\uline{1\textsuperscript{er} Cas} $Z$ est un fibré vectoriel $V(F)$, avec $F$ cohérent sur $X$. Alors le th. 12 nous fournit un $P$ cohérent sur $S$, tel que
\[
 \underline{\operatorname{Hom}}(F,f^*(M))\simeq\underline{\operatorname{Hom}}(P,M),
\]
et on voit tout de suite que $f_*(V(F))\simeq V(P)$.

\uline{2\textsuperscript{\unsure{ième|ème}} Cas} Cas Général. On sait qu'on a
% archive sheet number: 22
